\scanpage{196}

% B01 / text / top, proposition and additions over first lines
\deleted{Considérons}
\uline{Proposition}\quad Soit $E$ un torseur sous un
groupe $\Lambda_E$ \inserted{$\subset\unsure{\mathfrak S|\mathfrak G}_E$} cyclique d'ordre $n$. \inserted{$(n\in\mathbb N,\ \deleted{\unsure{3|2}}\ n\geq\unsure{2|3})$} Considérons
l'ens. $\overrightarrow C(E)$ des structures de contour
\deleted{condition}\inserted{\uline{orienté}} sur $E$, invariées par $\Lambda_E$.

% B02 / text / upper middle, bijection and first marginal note
\deleted{On a} \deleted{Pour tout} \deleted{Pour tout}
Cet ens est en correspondance biunivoque
avec $\Lambda_E$, en associant à tout $u\in\Lambda_E$
l'unique structure de contour orienté, pour
laquelle $u$ soit l'opération de ``rotation \inserted{génératrice}'' associée à
l'orientation.
% slanted
\margin{$\Lambda_E$ est simplement transitif sur l'ens des repères directs de chacune des structures envisagées (à l'exception des deux cas \unsure{$\mathrm{tot}^{\mathrm t}$|tout à fait} dégénérés, en décomposant en 1-gônes ou bigônes --)}

% B03 / text / middle, proof and NB in the left margin
En effet, la donnée d'une structure $\overrightarrow C$
de contour orienté, sur un ens fini \inserted{$E$}, revient
à celle d'un $u\in\mathfrak S_E$ -- la rotation \unsure{générique|génératrice} du contour orienté. Si un groupe
$\Lambda_E$ opère sur $E$, dire que $\overrightarrow C$ est invariante
par ce groupe, signifie que \inserted{$u=u_{\overrightarrow C}$} l'est, i.e.
que $u$ commute à ce groupe. Or dans le
cas d'une structure de torseur, le
commutant est réduit à $\Lambda_E$ -- cqfd.
% slanted
\margin{\uline{NB} Ici on n'exclut pas les contours ``\unsure{ayant}'' des orbites d'ordre 1 et 2 -- ce qui correspond à une définition \unsure{intuitible|intuitive|intraitable|inutilement restrictive} dans l'optique où la structure est portée par l'ens des sommets \ldots Si on voulait l'éviter, il faudrait dans la proposition se borner aux $u\in\Lambda_E$ qui ne sont pas d'ordre 1 ou 2, ce qui revient à \inserted{éliminer} juste les \deleted{\unsure{sommets}} éléments du \unsure{ss-groupe|s-groupe} d'ordre 2 \inserted{$ {}_2\Lambda_E$}}

% B04 / text / foot, corollary 1
\uline{Corollaire 1} Les structures $C$ envisagées forment
un module libre de rg 1 sur $\Lambda_n=\mathbb Z/n\mathbb Z$
(savoir $\Lambda_E$ !). Les générateurs de ce
module correspondent aux structures
polygônales orientées sur $E$, invariées par $\Lambda_E$,
qui forment donc un \uline{torseur} \inserted{$\Lambda_E^*$} sur $(\mathbb Z/n\mathbb Z)^*$.

% archive sheet number: 195
